\documentclass[12pt,a4paper]{article}
% Preámbulo común de los reportes RE_x-y (un documento por resultado esperado).
% Cada RE_*.tex lo carga con % Preámbulo común de los reportes RE_x-y (un documento por resultado esperado).
% Cada RE_*.tex lo carga con % Preámbulo común de los reportes RE_x-y (un documento por resultado esperado).
% Cada RE_*.tex lo carga con \input{preambulo} justo después de \documentclass.

% === Codificación e idioma =========================================
\usepackage[utf8]{inputenc}
\usepackage[T1]{fontenc}
\usepackage{lmodern}
\usepackage[spanish,es-tabla,es-noquoting]{babel}
\usepackage{microtype}
\usepackage{setspace}
\usepackage{etoolbox}

% === Figuras, tablas y listas =======================================
\usepackage{graphicx}
\usepackage{amsmath}
\usepackage{amssymb}
\usepackage{array}
\usepackage{booktabs}
\usepackage{tabularx}
\usepackage{makecell}
\usepackage{threeparttable}
\usepackage{enumitem}
\usepackage[font=small,labelfont=bf,skip=8pt]{caption}
\usepackage{subcaption}
\usepackage{titlesec}
\usepackage{xcolor}
\usepackage{listings}
\usepackage{fancyhdr}

% === Esquemas y gráficos en LaTeX ===================================
\usepackage{tikz}
\usetikzlibrary{arrows.meta,positioning,calc,shapes.geometric,fit,backgrounds,
                decorations.pathreplacing,matrix,chains,quotes}
\usepackage{pgfplots}
\pgfplotsset{compat=1.18}
\usepgfplotslibrary{groupplots,statistics}

% === Geometría ======================================================
\usepackage{geometry}
\geometry{left=3cm, right=2.5cm, top=2.5cm, bottom=2.5cm}

% === Citas y enlaces (hyperref al final) ============================
\usepackage[round,authoryear]{natbib}
\usepackage[hidelinks,hypertexnames=false]{hyperref}
\hypersetup{pdfauthor={Fernando Nelson Candia Aroni}, pdflang={es}}

% ===================================================================
% FORMATO
% ===================================================================
\setstretch{1.2}
\AtBeginEnvironment{tabular}{\setstretch{1}}
\AtBeginEnvironment{tabularx}{\setstretch{1}}
\AtBeginEnvironment{thebibliography}{\setstretch{1}}
\captionsetup{font={small,stretch=1}}

\titleformat*{\section}{\setstretch{1}\raggedright\normalfont\Large\bfseries}
\titleformat*{\subsection}{\setstretch{1}\raggedright\normalfont\large\bfseries}
\titleformat*{\subsubsection}{\setstretch{1}\raggedright\normalfont\normalsize\bfseries}
\titleformat{\paragraph}[hang]{\setstretch{1}\normalfont\bfseries}{}{0pt}{}
\titlespacing*{\paragraph}{0pt}{2.6ex plus .6ex minus .2ex}{1.2ex}

% Etiqueta corta incrustada al inicio del párrafo.
\newcommand{\rotulo}[1]{\par\medskip\noindent\textbf{#1}\hskip 1em\ignorespaces}

\setlength{\emergencystretch}{2em}
\widowpenalty=10000
\clubpenalty=10000
\setlist[description]{style=unboxed, leftmargin=1.6em, labelindent=0pt,
                      topsep=1ex, itemsep=1.1ex, parsep=0pt}
\setlist[itemize]{topsep=1ex, itemsep=0.5ex, parsep=0pt}
\setlist[enumerate]{topsep=1ex, itemsep=0.5ex, parsep=0pt}

\newcolumntype{L}[1]{>{\raggedright\arraybackslash}p{#1}}
\newcolumntype{R}[1]{>{\raggedleft\arraybackslash}p{#1}}
\newcolumntype{C}[1]{>{\centering\arraybackslash}p{#1}}
\renewcommand{\arraystretch}{1.15}
\renewcommand{\theadfont}{\bfseries}
\renewcommand{\theadalign}{br}

% Código y comandos: monoespaciado pequeño, sin resaltar.
\lstset{
    basicstyle=\ttfamily\small,
    breaklines=true,
    columns=fullflexible,
    keepspaces=true,
    frame=leftline,
    framerule=0.6pt,
    rulecolor=\color{black!35},
    xleftmargin=1em,
    aboveskip=1ex,
    belowskip=1ex,
    literate={á}{{\'a}}1 {é}{{\'e}}1 {í}{{\'i}}1 {ó}{{\'o}}1 {ú}{{\'u}}1 {ñ}{{\~n}}1
             {→}{{$\to$}}1 {≥}{{$\geq$}}1 {·}{{$\cdot$}}1,
}
\renewcommand{\lstlistingname}{Listado}

% === Paleta, estilos TikZ, \code, \file, \sep, \todocita y \todo ===========
\input{estilos}

% === Encabezado del reporte =========================================
% \reporte{RE1.1}{OE1}{Título}: título, subtítulo (RE y OE) y autor.
\newcommand{\reporte}[3]{%
    \pagestyle{fancy}\fancyhf{}%
    \fancyhead[L]{\small #2\sep#1}\fancyhead[R]{\small\thepage}%
    \renewcommand{\headrulewidth}{0.4pt}%
    \begin{center}
        {\LARGE\bfseries #3}\\[1.2ex]
        {\large Resultado esperado #1 del objetivo específico #2}\\[1.2ex]
        {\normalsize Fernando Nelson Candia Aroni}
    \end{center}
    \vspace{1ex}\hrule\vspace{2ex}%
}
 justo después de \documentclass.

% === Codificación e idioma =========================================
\usepackage[utf8]{inputenc}
\usepackage[T1]{fontenc}
\usepackage{lmodern}
\usepackage[spanish,es-tabla,es-noquoting]{babel}
\usepackage{microtype}
\usepackage{setspace}
\usepackage{etoolbox}

% === Figuras, tablas y listas =======================================
\usepackage{graphicx}
\usepackage{amsmath}
\usepackage{amssymb}
\usepackage{array}
\usepackage{booktabs}
\usepackage{tabularx}
\usepackage{makecell}
\usepackage{threeparttable}
\usepackage{enumitem}
\usepackage[font=small,labelfont=bf,skip=8pt]{caption}
\usepackage{subcaption}
\usepackage{titlesec}
\usepackage{xcolor}
\usepackage{listings}
\usepackage{fancyhdr}

% === Esquemas y gráficos en LaTeX ===================================
\usepackage{tikz}
\usetikzlibrary{arrows.meta,positioning,calc,shapes.geometric,fit,backgrounds,
                decorations.pathreplacing,matrix,chains,quotes}
\usepackage{pgfplots}
\pgfplotsset{compat=1.18}
\usepgfplotslibrary{groupplots,statistics}

% === Geometría ======================================================
\usepackage{geometry}
\geometry{left=3cm, right=2.5cm, top=2.5cm, bottom=2.5cm}

% === Citas y enlaces (hyperref al final) ============================
\usepackage[round,authoryear]{natbib}
\usepackage[hidelinks,hypertexnames=false]{hyperref}
\hypersetup{pdfauthor={Fernando Nelson Candia Aroni}, pdflang={es}}

% ===================================================================
% FORMATO
% ===================================================================
\setstretch{1.2}
\AtBeginEnvironment{tabular}{\setstretch{1}}
\AtBeginEnvironment{tabularx}{\setstretch{1}}
\AtBeginEnvironment{thebibliography}{\setstretch{1}}
\captionsetup{font={small,stretch=1}}

\titleformat*{\section}{\setstretch{1}\raggedright\normalfont\Large\bfseries}
\titleformat*{\subsection}{\setstretch{1}\raggedright\normalfont\large\bfseries}
\titleformat*{\subsubsection}{\setstretch{1}\raggedright\normalfont\normalsize\bfseries}
\titleformat{\paragraph}[hang]{\setstretch{1}\normalfont\bfseries}{}{0pt}{}
\titlespacing*{\paragraph}{0pt}{2.6ex plus .6ex minus .2ex}{1.2ex}

% Etiqueta corta incrustada al inicio del párrafo.
\newcommand{\rotulo}[1]{\par\medskip\noindent\textbf{#1}\hskip 1em\ignorespaces}

\setlength{\emergencystretch}{2em}
\widowpenalty=10000
\clubpenalty=10000
\setlist[description]{style=unboxed, leftmargin=1.6em, labelindent=0pt,
                      topsep=1ex, itemsep=1.1ex, parsep=0pt}
\setlist[itemize]{topsep=1ex, itemsep=0.5ex, parsep=0pt}
\setlist[enumerate]{topsep=1ex, itemsep=0.5ex, parsep=0pt}

\newcolumntype{L}[1]{>{\raggedright\arraybackslash}p{#1}}
\newcolumntype{R}[1]{>{\raggedleft\arraybackslash}p{#1}}
\newcolumntype{C}[1]{>{\centering\arraybackslash}p{#1}}
\renewcommand{\arraystretch}{1.15}
\renewcommand{\theadfont}{\bfseries}
\renewcommand{\theadalign}{br}

% Código y comandos: monoespaciado pequeño, sin resaltar.
\lstset{
    basicstyle=\ttfamily\small,
    breaklines=true,
    columns=fullflexible,
    keepspaces=true,
    frame=leftline,
    framerule=0.6pt,
    rulecolor=\color{black!35},
    xleftmargin=1em,
    aboveskip=1ex,
    belowskip=1ex,
    literate={á}{{\'a}}1 {é}{{\'e}}1 {í}{{\'i}}1 {ó}{{\'o}}1 {ú}{{\'u}}1 {ñ}{{\~n}}1
             {→}{{$\to$}}1 {≥}{{$\geq$}}1 {·}{{$\cdot$}}1,
}
\renewcommand{\lstlistingname}{Listado}

% === Paleta, estilos TikZ, \code, \file, \sep, \todocita y \todo ===========
% Estilos compartidos por los reportes RE_x-y y las presentaciones: paleta, estilos TikZ,
% macros tipográficas y marcas de trabajo pendiente. Requiere xcolor y tikz ya cargados.

\newcommand{\code}[1]{\texttt{#1}}
\newcommand{\file}[1]{\texttt{#1}}

% === Paleta de los esquemas ========================================
\definecolor{azul}{HTML}{1E40AF}
\definecolor{celeste}{HTML}{00B4D8}
\definecolor{verde}{HTML}{22C55E}
\definecolor{rosa}{HTML}{FF5FA2}
\definecolor{ambar}{HTML}{F59E0B}
\definecolor{gris}{HTML}{6B7280}
\tikzset{
    caja/.style={draw=black!60, rounded corners=2pt, fill=white, align=center,
                 minimum height=7mm, inner sep=4pt, font=\small},
    etapa/.style={caja, fill=azul!8, draw=azul!70},
    dato/.style={caja, fill=black!4, draw=black!50, font=\small\ttfamily},
    decision/.style={caja, fill=ambar!15, draw=ambar!80!black},
    descarte/.style={caja, fill=rosa!12, draw=rosa!70!black, font=\footnotesize},
    flecha/.style={-{Stealth[length=2.2mm]}, thick, draw=black!65},
    nota/.style={font=\footnotesize, text=black!70, align=left},
}

% Separador tipográfico (el «·» directo no existe en T1).
\newcommand{\sep}{\,\textperiodcentered\,}

% === Marcas de trabajo pendiente ===================================
% \todocita{...}: dónde hace falta una referencia y para qué. Visibles en el PDF.
\newcommand{\todocita}[1]{\textcolor{red!75!black}{\small[\textsc{cita}: #1]}}
\newcommand{\todo}[1]{\textcolor{orange!80!black}{\small[\textsc{todo}: #1]}}


% === Encabezado del reporte =========================================
% \reporte{RE1.1}{OE1}{Título}: título, subtítulo (RE y OE) y autor.
\newcommand{\reporte}[3]{%
    \pagestyle{fancy}\fancyhf{}%
    \fancyhead[L]{\small #2\sep#1}\fancyhead[R]{\small\thepage}%
    \renewcommand{\headrulewidth}{0.4pt}%
    \begin{center}
        {\LARGE\bfseries #3}\\[1.2ex]
        {\large Resultado esperado #1 del objetivo específico #2}\\[1.2ex]
        {\normalsize Fernando Nelson Candia Aroni}
    \end{center}
    \vspace{1ex}\hrule\vspace{2ex}%
}
 justo después de \documentclass.

% === Codificación e idioma =========================================
\usepackage[utf8]{inputenc}
\usepackage[T1]{fontenc}
\usepackage{lmodern}
\usepackage[spanish,es-tabla,es-noquoting]{babel}
\usepackage{microtype}
\usepackage{setspace}
\usepackage{etoolbox}

% === Figuras, tablas y listas =======================================
\usepackage{graphicx}
\usepackage{amsmath}
\usepackage{amssymb}
\usepackage{array}
\usepackage{booktabs}
\usepackage{tabularx}
\usepackage{makecell}
\usepackage{threeparttable}
\usepackage{enumitem}
\usepackage[font=small,labelfont=bf,skip=8pt]{caption}
\usepackage{subcaption}
\usepackage{titlesec}
\usepackage{xcolor}
\usepackage{listings}
\usepackage{fancyhdr}

% === Esquemas y gráficos en LaTeX ===================================
\usepackage{tikz}
\usetikzlibrary{arrows.meta,positioning,calc,shapes.geometric,fit,backgrounds,
                decorations.pathreplacing,matrix,chains,quotes}
\usepackage{pgfplots}
\pgfplotsset{compat=1.18}
\usepgfplotslibrary{groupplots,statistics}

% === Geometría ======================================================
\usepackage{geometry}
\geometry{left=3cm, right=2.5cm, top=2.5cm, bottom=2.5cm}

% === Citas y enlaces (hyperref al final) ============================
\usepackage[round,authoryear]{natbib}
\usepackage[hidelinks,hypertexnames=false]{hyperref}
\hypersetup{pdfauthor={Fernando Nelson Candia Aroni}, pdflang={es}}

% ===================================================================
% FORMATO
% ===================================================================
\setstretch{1.2}
\AtBeginEnvironment{tabular}{\setstretch{1}}
\AtBeginEnvironment{tabularx}{\setstretch{1}}
\AtBeginEnvironment{thebibliography}{\setstretch{1}}
\captionsetup{font={small,stretch=1}}

\titleformat*{\section}{\setstretch{1}\raggedright\normalfont\Large\bfseries}
\titleformat*{\subsection}{\setstretch{1}\raggedright\normalfont\large\bfseries}
\titleformat*{\subsubsection}{\setstretch{1}\raggedright\normalfont\normalsize\bfseries}
\titleformat{\paragraph}[hang]{\setstretch{1}\normalfont\bfseries}{}{0pt}{}
\titlespacing*{\paragraph}{0pt}{2.6ex plus .6ex minus .2ex}{1.2ex}

% Etiqueta corta incrustada al inicio del párrafo.
\newcommand{\rotulo}[1]{\par\medskip\noindent\textbf{#1}\hskip 1em\ignorespaces}

\setlength{\emergencystretch}{2em}
\widowpenalty=10000
\clubpenalty=10000
\setlist[description]{style=unboxed, leftmargin=1.6em, labelindent=0pt,
                      topsep=1ex, itemsep=1.1ex, parsep=0pt}
\setlist[itemize]{topsep=1ex, itemsep=0.5ex, parsep=0pt}
\setlist[enumerate]{topsep=1ex, itemsep=0.5ex, parsep=0pt}

\newcolumntype{L}[1]{>{\raggedright\arraybackslash}p{#1}}
\newcolumntype{R}[1]{>{\raggedleft\arraybackslash}p{#1}}
\newcolumntype{C}[1]{>{\centering\arraybackslash}p{#1}}
\renewcommand{\arraystretch}{1.15}
\renewcommand{\theadfont}{\bfseries}
\renewcommand{\theadalign}{br}

% Código y comandos: monoespaciado pequeño, sin resaltar.
\lstset{
    basicstyle=\ttfamily\small,
    breaklines=true,
    columns=fullflexible,
    keepspaces=true,
    frame=leftline,
    framerule=0.6pt,
    rulecolor=\color{black!35},
    xleftmargin=1em,
    aboveskip=1ex,
    belowskip=1ex,
    literate={á}{{\'a}}1 {é}{{\'e}}1 {í}{{\'i}}1 {ó}{{\'o}}1 {ú}{{\'u}}1 {ñ}{{\~n}}1
             {→}{{$\to$}}1 {≥}{{$\geq$}}1 {·}{{$\cdot$}}1,
}
\renewcommand{\lstlistingname}{Listado}

% === Paleta, estilos TikZ, \code, \file, \sep, \todocita y \todo ===========
% Estilos compartidos por los reportes RE_x-y y las presentaciones: paleta, estilos TikZ,
% macros tipográficas y marcas de trabajo pendiente. Requiere xcolor y tikz ya cargados.

\newcommand{\code}[1]{\texttt{#1}}
\newcommand{\file}[1]{\texttt{#1}}

% === Paleta de los esquemas ========================================
\definecolor{azul}{HTML}{1E40AF}
\definecolor{celeste}{HTML}{00B4D8}
\definecolor{verde}{HTML}{22C55E}
\definecolor{rosa}{HTML}{FF5FA2}
\definecolor{ambar}{HTML}{F59E0B}
\definecolor{gris}{HTML}{6B7280}
\tikzset{
    caja/.style={draw=black!60, rounded corners=2pt, fill=white, align=center,
                 minimum height=7mm, inner sep=4pt, font=\small},
    etapa/.style={caja, fill=azul!8, draw=azul!70},
    dato/.style={caja, fill=black!4, draw=black!50, font=\small\ttfamily},
    decision/.style={caja, fill=ambar!15, draw=ambar!80!black},
    descarte/.style={caja, fill=rosa!12, draw=rosa!70!black, font=\footnotesize},
    flecha/.style={-{Stealth[length=2.2mm]}, thick, draw=black!65},
    nota/.style={font=\footnotesize, text=black!70, align=left},
}

% Separador tipográfico (el «·» directo no existe en T1).
\newcommand{\sep}{\,\textperiodcentered\,}

% === Marcas de trabajo pendiente ===================================
% \todocita{...}: dónde hace falta una referencia y para qué. Visibles en el PDF.
\newcommand{\todocita}[1]{\textcolor{red!75!black}{\small[\textsc{cita}: #1]}}
\newcommand{\todo}[1]{\textcolor{orange!80!black}{\small[\textsc{todo}: #1]}}


% === Encabezado del reporte =========================================
% \reporte{RE1.1}{OE1}{Título}: título, subtítulo (RE y OE) y autor.
\newcommand{\reporte}[3]{%
    \pagestyle{fancy}\fancyhf{}%
    \fancyhead[L]{\small #2\sep#1}\fancyhead[R]{\small\thepage}%
    \renewcommand{\headrulewidth}{0.4pt}%
    \begin{center}
        {\LARGE\bfseries #3}\\[1.2ex]
        {\large Resultado esperado #1 del objetivo específico #2}\\[1.2ex]
        {\normalsize Fernando Nelson Candia Aroni}
    \end{center}
    \vspace{1ex}\hrule\vspace{2ex}%
}

\hypersetup{pdftitle={RE2.5 Código y modelo disponibles en un repositorio}}

\begin{document}
\reporte{RE2.5}{OE2}{Código y modelo disponibles en un repositorio}

\section{Alcance}

\rotulo{Medio de verificación.} Repositorio público
\url{https://github.com/nhrot-fc/tesis-primate} y sus \emph{releases}. Este documento es un
índice y no tiene cuaderno.

\rotulo{Indicador.} La URL del repositorio permite localizar el código de curación,
entrenamiento, evaluación e inferencia, las instrucciones de ejecución y los artefactos del
modelo final publicados (Tabla~\ref{tab:donde}).

\section{Repositorio}

\rotulo{URL.} \url{https://github.com/nhrot-fc/tesis-primate}, público, con historia completa
(más de cien \emph{commits}) y \emph{releases} etiquetados por versión
(\code{v0.1.0}, \code{v0.1.5}). La versión del paquete es la de \file{pyproject.toml}.

\rotulo{Estructura.} La Figura~\ref{fig:arbol} muestra el árbol de primer nivel y qué
resultado esperado documenta cada rama; la Tabla~\ref{tab:donde} lleva de cada verbo del
indicador al archivo que lo implementa.

\begin{figure}[htbp]
    \centering
    \begin{minipage}{0.92\textwidth}
\begin{lstlisting}
tesis-primate/
  README.md         instalacion del paquete portable y uso del visor
  pyproject.toml    dependencias, extras y grupos (uv)
  evaluation.sh     volcado + comparacion de todas las corridas
  src/              el pipeline entero (RE2.2, mapa del codigo)
  notebooks/        auditoria, comparacion visual, figuras, aves
  research/         la tesis (main.tex), figuras, reportes RE_*
  deploy/           paquete portable para Windows
  docs/             protocolo de comparacion, requisitos, notas
  data/   (no versionado)  raw/, cleaned/, processed/, yolo/
  runs/   (no versionado)  corridas: checkpoints, volcados
  hf/     (no versionado)  backbones preentrenados
\end{lstlisting}
    \end{minipage}
    \caption{Árbol del repositorio}
    \label{fig:arbol}
\end{figure}

\begin{table}[htbp]
    \centering
    \caption{Ubicación de cada elemento del indicador}
    \label{tab:donde}
    \small
    \begin{tabular}{@{}L{3.2cm}L{6.6cm}L{4.0cm}@{}}
        \toprule
        \textbf{Qué} & \textbf{Dónde} & \textbf{Documento} \\
        \midrule
        Curación & \file{src/prepare\_annotations.py}, \file{src/data/annotations.py}, \file{src/data/species.py} & RE1.1 \\
        Conjunto y partición & \file{src/prepare\_data.py}, \file{src/data/manifest.py}, \file{src/data/cache.py} & RE1.2, RE1.3 \\
        Arquitecturas & \file{src/models/} & RE2.1 \\
        Entrenamiento & \file{src/train.py}, \file{src/training/} & RE2.2 \\
        Evaluación & \file{src/dump\_predictions.py}, \file{src/compare\_models.py}, \file{src/evaluation/}, \file{evaluation.sh} & RE2.1 (protocolo), RE2.3 \\
        Inferencia & \file{src/detect.py}, \file{src/inference/}, \file{src/main.py} y \file{src/viewer/} & RE2.4 \\
        Instrucciones & \file{README.md}, \file{docs/system\_requirements.md}, \file{pyproject.toml} & RE2.2 \\
        Modelos publicados & \emph{Releases}: \code{detector-<versión>-\allowbreak model-<nombre>.zip} & esta sección \\
        Paquete portable & \emph{Releases}: \code{detector-<versión>-\allowbreak win64-\{cpu,cuda\}.zip}; \file{deploy/} & esta sección \\
        \bottomrule
    \end{tabular}
\end{table}

\section{Contenido publicado}

\subsection{Código}

Todo el pipeline, tal como lo describe RE2.2, más los cuadernos que generan las figuras de la
tesis y de estos reportes. El código corre con \code{uv sync} y los comandos de RE2.2; no
depende de nada que no esté declarado en \file{pyproject.toml} y \file{uv.lock}.

\subsection{Modelos}

Cada \emph{release} lleva un zip por modelo con la carpeta \code{models/<nombre>/} (el
\emph{checkpoint} autocontenido de RE2.2 más \file{operating\_point.json}) y, cuando el
modelo lo necesita, la copia de su \emph{backbone} preentrenado en \code{hf/}. El zip se
descomprime dentro de la carpeta del paquete portable, o de un clon del repositorio, y el
modelo aparece en el selector del visor y en \file{detect.py} sin ningún paso más. La
Tabla~\ref{tab:release} lista los artefactos del \emph{release} actual.

\begin{table}[htbp]
    \centering
    \begin{threeparttable}
        \caption{Artefactos del \emph{release} \code{v0.1.5}}
        \label{tab:release}
        \small
        \begin{tabular}{@{}L{6.6cm}rL{4.8cm}@{}}
            \toprule
            \textbf{Archivo} & \thead{Tamaño} & \textbf{Contenido} \\
            \midrule
            \file{detector-0.1.5-win64-cpu.zip} & 0,5\,GB & Programa completo para Windows de 64 bits, sólo CPU \\
            \file{detector-0.1.5-win64-\allowbreak cuda.zip.001/.002} & 3,1\,GB & Lo mismo con soporte NVIDIA; \file{join.bat} une las partes \\
            \file{detector-0.1.5-model-frcnn.zip} & 161\,MB & Faster R-CNN R50-FPN v2 \\
            \file{detector-0.1.5-model-\allowbreak yolo26s\_coco.zip} & 24\,MB & YOLO26s \\
            \file{detector-0.1.5-model-\allowbreak detr\_t10\_pcen.zip} & 652\,MB & AST-Deformable-DETR con PCEN, con su AST \\
            \file{SHA256SUMS.txt} & --- & Sumas de verificación de todo lo anterior \\
            \bottomrule
        \end{tabular}
        \begin{tablenotes}[flushleft]
            \footnotesize
            \item Los tres modelos del \emph{release} se entrenaron sobre la versión anterior
            del caché (antes de unir las copias de grabaciones, RE1.1). \todo{Publicar
            \code{yolo26s\_v2}, el modelo final de RE2.3 y RE2.4, como
            \file{detector-0.1.6-model-yolo26s\_v2.zip}, y retirar o marcar como anteriores los
            tres zips actuales.}
        \end{tablenotes}
    \end{threeparttable}
\end{table}

\subsection{Paquete portable}

El paquete de Windows (\file{deploy/build\_windows.py}) empaqueta un Python embebido con las
dependencias resueltas por \code{uv} para Windows, el código de \file{src/} y dos lanzadores
compilados desde \file{deploy/launcher/launcher.c}: \file{Detector.exe} abre el visor (vistas
\emph{Spectrogram} y \emph{Batch}) y \file{detect.exe} es la línea de comandos de
\file{detect.py}. No exige instalar Python, CUDA ni permisos de administrador; el usuario
descomprime, añade un modelo y arrastra un audio o una carpeta. \file{deploy/release.sh} sube
los zips a GitHub con sus sumas SHA-256. \todo{El paquete se construyó y probó en Linux con
Wine ausente: falta una prueba en una máquina Windows real antes de publicar el \emph{release}
como definitivo (hoy está en borrador).}

\section{Contenido no publicado}

\begin{itemize}
    \item \textbf{Las grabaciones y las anotaciones primarias} (\file{data/raw/}) pertenecen
          al equipo de investigación; el repositorio no las contiene ni las referencia por
          ubicación. RE1.3 deja abierta la decisión de liberar el conjunto derivado.
    \item \textbf{Las corridas} (\file{runs/}): sus \emph{checkpoints} y volcados suman
          decenas de gigabytes; se publican los \emph{checkpoints} elegidos, no las corridas.
    \item \textbf{Los pesos preentrenados de terceros} (AST de AudioSet, Deformable DETR de
          COCO) se descargan de Hugging Face en el primer uso y viajan dentro del zip del
          modelo que los necesita, con la licencia de su autor.
\end{itemize}

\section{Comandos de reproducción}

\begin{lstlisting}[language=bash]
git clone https://github.com/nhrot-fc/tesis-primate.git && cd tesis-primate
uv sync --no-dev --extra yolo                       # o --extra detr para los DETR
unzip detector-0.1.5-model-yolo26s_coco.zip          # deja models/yolo26s_coco/
uv run python src/detect.py grabaciones/            # una tabla de Raven por audio
uv run python src/main.py grabaciones/alguna.wav    # o el visor
\end{lstlisting}

Sin GPU los cinco comandos funcionan igual; en Windows, el paquete portable reemplaza los dos
primeros por descomprimir un zip.

\section{Licencia y condiciones de uso}

\todo{El repositorio no tiene archivo \file{LICENSE}. Decidir y añadir: para el código, una
licencia permisiva (MIT o Apache-2.0) es lo habitual en software de investigación; para los
pesos del modelo, declarar que derivan de AudioSet/COCO y heredan las condiciones de sus
\emph{backbones}; para los datos, remitir a la decisión del equipo (RE1.3). Sin licencia, el
código es ``todos los derechos reservados'' por defecto y nadie puede reutilizarlo
legalmente.}

\todocita{principios de citación de software (Smith et al. 2016) y, si se deposita una
versión en Zenodo con DOI, la guía de GitHub--Zenodo; también la práctica de publicar pesos
con \emph{model cards} (Mitchell et al. 2019), que la ficha de RE1.3 y el
\file{operating\_point.json} cubren parcialmente}

\section{Versionado y trazabilidad}

Cada \emph{checkpoint} lleva dentro la descripción del conjunto con el que se entrenó (RE2.2)
y cada \emph{release} lleva sumas SHA-256. Las cifras de los reportes RE1.x y RE2.x se
generan con los cuadernos homónimos de \file{research/reportes/} sobre los artefactos de
\file{runs/} y \file{data/processed/}, de modo que una cifra se rastrea a un \emph{commit}, un
\file{meta.json} y un \file{best.pt}. \todo{Etiquetar el \emph{commit} con el que se
generaron estos reportes (\code{git tag reportes-oe1-oe2}) y depositar una instantánea en
Zenodo para obtener un DOI citable en la tesis.}

\clearpage
\bibliographystyle{plainnat}
\bibliography{../references}
\end{document}
